\section{Selection for logarithmic models on products}
\label{sec:logarithmic-model-selection}

For a logarithmic reference with a positive residual background, removing the divisor turns the singular exponential equation into a bounded equation. On a product with a curve, the selected residual can then be determined by the Green operator. This also controls simultaneous changes of the reference, background and density, for which the centering constant must follow the changing data.

Let $(Y,\kappa)$ be a connected compact Hermitian manifold of dimension $k\geq1$, with $\int_Y\kappa^k=1$, and put $X=\mathbb P^1\times Y$, so that $n=k+1$. Write $\omega_1$ for the mass-one Fubini--Study form on $\mathbb P^1$ and $\mathrm{d}W=\omega_1\wedge\kappa^k$; pullbacks to $X$ are implicit. Fix $a,d>0$, a smooth positive nonconstant function $b$ on $\mathbb P^1$, and a continuous positive function $H$ on $\mathbb P^1$. Set
\[
 \begin{aligned}
 \theta&=(a+d)\omega_1+b\kappa,&\eta&=a\omega_1+b\kappa,\\
 f&=n b^kH,&m&=\int_{\mathbb P^1}H\omega_1,\qquad c_0=\frac am.
 \end{aligned}
\]
Both forms are positive and nonclosed: the component $\mathrm{d}b\wedge\kappa$ of horizontal degree one cannot cancel the vertical component $b\,\mathrm{d}\kappa$. For homogeneous coordinates $[Z_0:Z_1]$, put
\[
 \ell=\frac{|Z_0|^2}{|Z_0|^2+|Z_1|^2},\qquad
 D=\{Z_0=0\}\times Y,\qquad \psi=d\log\ell.
\]
The Poincar\'e--Lelong formula gives $\theta+\ddc\psi=\eta+d[D]$.

\subsection{The ordinary model and the greatest solutions}

Put $K=a+d$, $R=a/d$ and $C_*=K\log(K/d)-a\log(a/d)$. In the coordinate $z=Z_1/Z_0$ define
\begin{equation}\label{eq:logarithmic-product-model}
 \phi_{a,d}=\begin{cases}
 0,&|z|^2\leq R,\\
 a\log|z|^2+C_*-K\log(1+|z|^2),&|z|^2>R,
 \end{cases}
 \qquad \phi_{a,d}|_D=-\infty.
\end{equation}

\begin{proposition}\label{prop:logarithmic-product-model}\index{Model potential!logarithmic product}
The function $\phi_{a,d}$ is an ordinary model reference for $\theta$, with
\[
 P_\theta[\psi]=P_\theta[\phi_{a,d}]=\phi_{a,d},\qquad
 \psi\leq\phi_{a,d}\leq\psi+C_*,\qquad \sup_X\phi_{a,d}=0.
\]
For $d/(a+d)<\delta<1$, the envelope $P_{\delta\theta}(\phi_{a,d})$ is not identically $-\infty$. The model has finite positive Monge--Amp\`ere measure
\[
 \mathrm{MA}_\theta(\phi_{a,d})
       =n b^kK\mathbf1_{\{|z|^2<R\}}\,\mathrm{d}W.
\]
The difference $\phi_{a,d}-\psi$ extends continuously across $D$. If $(a_j,d_j)\to(a,d)$ in $(0,+\infty)^2$, these residual extensions converge uniformly.
\end{proposition}
\begin{proof}
First take the ordinary envelope on the curve with background $K\omega_1$ and obstacle $\min(\psi+C,0)$. Rotation of $z$ preserves both the obstacle and the admissible family, so the envelope is rotation invariant. With $s=\log|z|^2$ and $F(s)=K\log(1+\mathrm{e}^s)$, its profile $u(s)$ has the property that $u+F$ is convex and lies below $F$. The comparison with $\psi+C$ bounds the final slope of $u+F$ by $a$. At $s_*=\log R$, one has $F'(s_*)=a$. Above $s_*$, convexity and the final slope bound give
\[
 u(s)+F(s)\leq F(s_*)+a(s-s_*),
\]
and below $s_*$ the obstacle gives $u\leq0$. The function in \eqref{eq:logarithmic-product-model} attains both bounds. Its convex profile, including its logarithmic value at infinity, gives a curve potential, and direct calculation gives $\psi\leq\phi_{a,d}\leq\psi+C_*$. It is therefore a competitor for every $C\geq C_*$, so the curve envelope is exactly $\phi_{a,d}$.

A product competitor restricts to each curve $\mathbb P^1\times\{y\}$ as either the constant $-\infty$ or a curve competitor. It is thus bounded above by $\phi_{a,d}$. Conversely, the curve model is a product competitor because the vertical form $b\kappa$ is positive. Hence the product envelope is the same function. Ordinary-type invariance of the singularity envelope gives its model property. Moreover,
\[
 \delta\theta+\ddc\psi
   =(\delta K-d)\omega_1+\delta b\kappa+d[D]\geq0,
\]
which proves the smaller-background assertion.

Off $D$, the horizontal curvature of $\phi_{a,d}$ is $K\mathbf1_{\{|z|^2<R\}}\omega_1$. The first derivatives of its convex profile agree at $R$, so there is no boundary-circle atom. Expanding the product off $D$ gives the displayed measure, and the non-pluripolar product assigns $D$ zero mass. Finally, in terms of $\ell$ the residual extension is
\[
 \phi_{a,d}-\psi=\begin{cases}
 -d\log\ell,&\ell\geq d/K,\\
 a\log(1-\ell)+C_*,&\ell<d/K.
 \end{cases}
\]
The formulas agree at the moving interface and extend to both endpoints. On compact positive parameter ranges each logarithm is used away from its singular endpoint. Joint continuity on the resulting compact set proves uniform variation.
\end{proof}

\begin{lemma}\label{lem:logarithmic-greatest-solutions}
For every $\sigma>0$, the equation
\begin{equation}\label{eq:logarithmic-product-exponential}
 \mathrm{MA}_\theta(s_\sigma)=\mathrm{e}^{\sigma s_\sigma}f\,\mathrm{d}W
\end{equation}
has a solution of ordinary type $\phi_{a,d}$, unique in that type. It is the greatest exponential subsolution among all $u\in\PSH(X,\theta)$ with $u\preceq\phi_{a,d}$, and its measure has finite mass.
\end{lemma}
\begin{proof}
The one-sided ordinary class has the description
\begin{equation}\label{eq:logarithmic-one-sided-removal}
 u\preceq\psi\quad\Longleftrightarrow\quad
 u=\psi+r,\quad r\in\PSH(X,\eta),\qquad
 \mathrm{MA}^{\mathrm{ext}}_\theta(\psi+r)
       =\mathrm{MA}^{\mathrm{ext}}_\eta(r).
\end{equation}
This is \cref{lem:one-sided-divisor-correspondence} for the pullback of the hyperplane bundle of $\mathbb P^1$ with its Fubini--Study metric, the section $Z_0$ and the weight $d$; the residual form is $\eta$, and $\sup_X\psi=0$. Indeed, $u\preceq\psi$ gives $\nu(u,D)\geq d$, and divisorial removal, \cref{lem:nonclosed-divisorial-removal}, shows that $\theta+\ddc u-d[D]\geq0$, so that $u-\psi$ has an $\eta$-plurisubharmonic representative $r$. The bounded ordinary type corresponds exactly to bounded $r$, as in \cref{thm:ordinary-divisor-normalization}.

The function $g_\sigma=\mathrm{e}^{\sigma\psi}f$ extends continuously by zero on $D$ and is positive elsewhere. Since $\eta>0$, \cref{thm:nguyen-exp} gives a continuous solution $R_\sigma$ of
\[
 \mathrm{MA}_\eta(R_\sigma)=\mathrm{e}^{\sigma R_\sigma}g_\sigma\,\mathrm{d}W.
\]
For this bounded density, the greatest solution of \cref{thm:individual-exponential-subsolution} is bounded and coincides with $R_\sigma$ by \cref{cor:bounded-exponential-comparison}. Thus it dominates every $\eta$-plurisubharmonic exponential subsolution. The correspondence \eqref{eq:logarithmic-one-sided-removal} transfers this order to the entire one-sided ordinary class and gives $s_\sigma=\psi+R_\sigma$. Bounded residual uniqueness gives uniqueness in the ordinary type. The bounded residual equation has finite mass, so the singular equation does too.
\end{proof}

\Needspace{12\baselineskip}
\subsection{The selected limit with varying data}

For $x=[x_0:x_1]$ and $z=[z_0:z_1]$ in $\mathbb P^1$, let
\[
 \operatorname{dist}_{\mathrm{chord}}(x,z)=\frac{|x_0z_1-x_1z_0|}{|x|\,|z|}
\]
be the chordal distance, normalized so that antipodal points have distance one, and put
\[
 G(x,z)=\log\operatorname{dist}_{\mathrm{chord}}(x,z)^2+1.
\]
For $z=[0:1]$ one has $G(x,z)=\log\ell(x)+1$. This kernel satisfies $\ddc_xG=\delta_z-\omega_1$ and $\int G(x,z)\omega_1(x)=0$. For a bounded density $q$, write $G*q(x)=\int G(x,z)q(z)\omega_1(z)$. Define
\begin{equation}\label{eq:logarithmic-selected-residual}
 \begin{aligned}
 A_0&=\frac{aH}{m},&w_0&=G*(A_0-a),&S_0&=\sup_X(\psi+w_0),\\
 \mathrm{d}\nu_0&=\frac Hm\,\mathrm{d}W,&J&=\int_X(\psi+w_0)\,\mathrm{d}\nu_0.
 \end{aligned}
\end{equation}
The function $w_0$ is continuous, has mean zero on the curve and satisfies $\ddc w_0=(A_0-a)\omega_1$. We use the probability $\nu_0$ to normalize the separated potentials.

\Needspace{14\baselineskip}
\begin{theorem}\label{thm:moving-logarithmic-selection}\index{Monge--Ampere equation@Monge--Amp\`ere equation!selection for logarithmic models}
Let $a_j,d_j>0$ converge to $a,d>0$, let smooth positive functions $b_j$ converge uniformly to $b$, and let continuous positive functions $H_j$ converge uniformly to $H$. For any $\sigma_j>0$ tending to zero, put
\[
 \begin{aligned}
 \theta_j&=(a_j+d_j)\omega_1+b_j\kappa,&\psi_j&=d_j\log\ell,&\phi_j&=\phi_{a_j,d_j},\\
 f_j&=n b_j^kH_j,&m_j&=\int_{\mathbb P^1}H_j\omega_1,&c_{0,j}&=a_j/m_j.
 \end{aligned}
\]
Let $s_j$ be the greatest solution of \eqref{eq:logarithmic-product-exponential} for these data, and set $v_j=s_j-\sup_Xs_j$, $c_j=\mathrm{e}^{\sigma_j\sup_Xs_j}$. Then
\[
 v_j\longrightarrow v_0=\psi+w_0-S_0,\qquad
 s_j-\sigma_j^{-1}\log c_{0,j}\longrightarrow s_0=\psi+w_0-J
\]
in $L^1(X)$, in capacity, and locally uniformly on $X\setminus D$. Moreover,
\begin{equation}\label{eq:logarithmic-scalar-error}
 c_j\longrightarrow c_0,\qquad
 \sigma_j^{-1}\log(c_j/c_{0,j})\longrightarrow S_0-J,
\end{equation}
and $\mathrm{MA}_{\theta_j}(v_j)\to c_0 f\,\mathrm{d}W$ in total variation. The limit has ordinary type $\phi_{a,d}$ and belongs to $\mathcal E_{\mathrm{env}}(X,\theta,\phi_{a,d})$, with
\[
 \mathrm{MA}_\theta(v_0)=c_0 f\,\mathrm{d}W.
\]
There is no restriction on the relative rates of parameter and exponent convergence.
\end{theorem}
\begin{proof}
We construct uniformly bounded residuals on the curve, then use their equation to determine both the normalized limit and the scalar error.

\emph{Uniform construction.} The kernel has uniformly integrable logarithmic tails and
$C_G=\sup_x\int|G(x,z)|\omega_1(z)<\infty$. These facts follow by rotation invariance and the radial integral of $|\log\ell+1|$ on $0<\ell<1$; the Poincar\'e--Lelong formula and $\int_0^1\log\ell\,\mathrm{d}\ell=-1$ give its curvature and mean. Truncating the kernel shows that convolution of bounded densities is continuous.

Temporarily omit subscripts. On continuous mean-zero functions on the curve define
\[
 T_\sigma(w)=G*(A(w)-a),\qquad
 A(w)=\frac{aH\mathrm{e}^{\sigma(\psi+w)}}{Z(w)},\qquad
 Z(w)=\int_{\mathbb P^1}H\mathrm{e}^{\sigma(\psi+w)}\omega_1.
\]
For all late parameters choose fixed bounds $a\leq a_{\max}$, $0<d_{\min}\leq d\leq d_{\max}$ and $0<h_{\min}\leq H\leq h_{\max}$. Put $B=4a_{\max}h_{\max}/h_{\min}$ and choose $R_0>C_G(B+a_{\max})$. On the ball $\|w\|_\infty\leq R_0$, for sufficiently small $\sigma$,
\[
 Z(w)\geq\frac{h_{\min}\mathrm{e}^{-\sigma R_0}}{1+\sigma d_{\max}}
       \geq h_{\min}/2,\qquad 0\leq A(w)\leq B.
\]
Here $\int\ell^{\sigma d}\omega_1=(1+\sigma d)^{-1}$. Thus $T_\sigma$ maps the mean-zero ball into itself. Since $\int A(w)\omega_1=a$, its density derivative in direction $h$ is
\[
 \sigma A(w)\left(h-a^{-1}\int hA(w)\omega_1\right).
\]
It follows that $\|T_\sigma(w)-T_\sigma(v)\|_\infty\leq2\sigma B C_G\|w-v\|_\infty$. The map is a contraction for all sufficiently small exponents, uniformly over the late parameters. Its fixed point $w_j$ has mean zero and $\|w_j\|_\infty\leq R_0$.

Write $Z_j=Z(w_j)$ and $A_j=A(w_j)$. The fixed-point equation gives $\ddc w_j=(A_j-a_j)\omega_1$, and the potential
\begin{equation}\label{eq:logarithmic-exponential-residual}
 \psi_j+w_j+\sigma_j^{-1}\log(a_j/Z_j)
\end{equation}
has residual form $A_j\omega_1+b_j\kappa$, continuous and semipositive. On coordinate balls, mollification of $w_j$ converges together with its continuous complex Hessian. The local bounded product, expanded after a common strict smooth correction, therefore equals $n b_j^kA_j\,\mathrm{d}W$. Removing $D$ identifies the singular product with the same measure. The identity
\[
 n b_j^kA_j=\mathrm{e}^{\sigma_j(\psi_j+w_j)}(a_j/Z_j)f_j
\]
shows that \eqref{eq:logarithmic-exponential-residual} solves the exponential equation in the ordinary type. By \cref{lem:logarithmic-greatest-solutions}, it is $s_j$.

\emph{Residual and normalized limits.} Uniform boundedness and dominated convergence give $Z_j\to m$ and $A_j\to A_0$ in $L^1(\omega_1)$, with a common bound on $A_j$. For uniformly bounded densities $q_j\to0$ in $L^1$, one has
\[
 \|G*q_j\|_\infty\leq\|q_j\|_\infty
       \sup_x\int_{\{|G(x,z)|>M\}}|G(x,z)|\omega_1(z)
       +M\|q_j\|_1.
\]
Applied to $q_j=A_j-a_j-A_0+a$, first letting $j\to\infty$ and then $M\to\infty$ proves $w_j\to w_0$ uniformly. Although $A_j=0$ on $D$ and $A_0>0$ there, the $L^1$ convergence suffices for this argument.

Put $S_j=\sup_X(\psi_j+w_j)$. A point with $\ell=1$ supplies a uniformly bounded competing value. Since $d_j\geq d_{\min}>0$, the upper bound $d_{\min}\log\ell+R_0$ tends uniformly to $-\infty$ near $D$, so all maximizing points stay in a fixed compact subset of $X\setminus D$. Uniform convergence there gives $S_j\to S_0$. Hence $v_j=\psi_j+w_j-S_j$ has the asserted $L^1$ and local uniform limits. For capacity convergence, the complements of $\{\psi<-T\}$ are compact subsets of $X\setminus D$. Choose a fixed Hermitian form for which $\psi$ is plurisubharmonic and apply \eqref{eq:global-capacity-tail}; their omitted sublevels have capacity at most $C/T$. Local uniform convergence, followed by $T\to\infty$, proves convergence in capacity, with the common poles omitted.

The measure formula and uniform convergence of $b_j$ give total-variation convergence to $n b^kA_0\,\mathrm{d}W=c_0 f\,\mathrm{d}W$. The same bounded product calculation gives this equation for $v_0$. The residual extensions in \cref{prop:logarithmic-product-model} show directly that $v_0\sim\phi_{a,d}$. Hence $v_0$ belongs to the relative envelope class $\mathcal E_{\mathrm{env}}(X,\theta,\phi_{a,d})$ of \cref{def:relative-envelope-class}, which contains the whole ordinary type of $\phi_{a,d}$ by \cref{prop:relative-elementary}.

\emph{Scalar centering.} The exact formula is $c_j=a_j\mathrm{e}^{\sigma_j S_j}/Z_j$. In the difference quotient
\[
 \frac{Z_j-m_j}{\sigma_j}
   =\int H_j\frac{\mathrm{e}^{\sigma_j(\psi_j+w_j)}-1}{\sigma_j}\omega_1,
\]
the negative part is bounded by $h_{\max}(d_{\max}|\log\ell|+R_0)$ and the positive part by a constant. Dominated convergence gives a limit $\int H(\psi+w_0)\omega_1=mJ$. Since $m_j\to m>0$,
\[
 \sigma_j^{-1}\log(Z_j/m_j)\longrightarrow J.
\]
The identity $c_j/c_{0,j}=m_j\mathrm{e}^{\sigma_j S_j}/Z_j$ proves \eqref{eq:logarithmic-scalar-error}. Subtracting $\sigma_j^{-1}\log c_{0,j}$ in \eqref{eq:logarithmic-exponential-residual} leaves $\psi_j+w_j-\sigma_j^{-1}\log(Z_j/m_j)$, proving the centered limit in all three topologies.
\end{proof}

\Needspace{27\baselineskip}
\subsection{The critical coefficient and separated order}

\Needspace{10\baselineskip}
\begin{corollary}\label{cor:logarithmic-critical-order}\index{Normalization scalar!centering for moving references}
For the fixed data, $c_0$ is the largest coefficient $c>0$ for which some $u\in\PSH(X,\theta)$ with $u\preceq\phi_{a,d}$ satisfies $\mathrm{MA}^{\mathrm{ext}}_\theta(u)\geq c f\,\mathrm{d}W$. Among separated potentials $u=\psi+r(x)\preceq\phi_{a,d}$, the critical subsolutions, those with $\mathrm{MA}^{\mathrm{ext}}_\theta(u)\geq c_0f\,\mathrm{d}W$, are exactly the translates of $\psi+w_0$. Thus $v_0$ is their unique supremum-zero member, and $s_0$ is their greatest member with $\int_Xu\,\mathrm{d}\nu_0\leq0$.
\end{corollary}
\begin{proof}
For a separated critical subsolution, \eqref{eq:logarithmic-one-sided-removal} and the product expansion give $\mathrm{MA}^{\mathrm{ext}}_{a\omega_1}(r)\geq A_0\omega_1$. Closed curve truncations have mass $a$, so the non-pluripolar measure has mass at most $a$. Since $\int A_0\omega_1=a$, it must equal $A_0\omega_1$. The positive closed current $a\omega_1+\ddc r$ also has mass $a$, so it has no remaining measure on the pole set. Consequently $\ddc r=(A_0-a)\omega_1$ as distributions, and $r=w_0+C$ because harmonic distributions on $\mathbb P^1$ are constant. Supremum zero gives $C=-S_0$, while the mean condition gives $C\leq-J$.

For a general $u\preceq\phi_{a,d}$ with coefficient $c>0$, write $u=\psi+r$ by \eqref{eq:logarithmic-one-sided-removal} and put $B_r=\sup_Xr$. Since $\psi\leq0$, the potential $u+\sigma^{-1}\log c-B_r$ is an exponential subsolution for every $\sigma>0$. Since $s_\sigma$ is the greatest exponential subsolution in the one-sided class, by \cref{lem:logarithmic-greatest-solutions},
\[
 u+\sigma^{-1}\log c-B_r\leq s_\sigma.
\]
At a point off $D$ where $r$ is finite, subtract $\sigma^{-1}\log c_0$, multiply by $\sigma$ and use the fixed-data case of \cref{thm:moving-logarithmic-selection}. This yields $\log c\leq\log c_0$. The potential $v_0$ attains $c_0$.
\end{proof}

The scalar in the centering follows the varying data. Replacing $c_{0,j}$ by $c_0$ adds $\sigma_j^{-1}\log(c_{0,j}/c_0)$ to the centered potential, and this term need not vanish. For example, take $a=1$, $d=1/4$, $H=1$ and $b=2+t$, where $t$ is the sphere height. Then $w_0=0$, $S_0=0$, $J=-1/4$ and $c_0=1$, so $s_0=\psi+1/4$ and the scalar-error limit is $1/4$. If $a_j=1+\sqrt{\sigma_j}$ and the other data are fixed, the appropriate centering uses $c_{0,j}=a_j$. Centering by $1$ instead adds
\[
 \frac{\log(1+\sqrt{\sigma_j})}{\sigma_j}\longrightarrow+\infty.
\]
These examples can be taken on a non-K\"ahler manifold by choosing $Y$ to be the Hopf surface with the unit-volume form $\kappa$ of \cref{ex:non-dini-hopf-product-adjoint}. The product cannot be K\"ahler, since a K\"ahler form would restrict to one on each slice $\{x\}\times Y$. The pole divisor remains present in the selected ordinary type.

\Needspace{12\baselineskip}
\subsection{Weak adjoints on a product surface}
\label{subsec:singular-product-adjoints}

On the product of two projective lines, the averaging probability of \eqref{eq:logarithmic-selected-residual} is the full supporting probability. The continuous positive residual defines a weak linearization, and spherical harmonics give invariant-measure uniqueness without a regularity assumption.

In this subsection take $Y=\mathbb P^1_y$, $\kappa=\omega_2$ and $k=1$. Keep $a,d,b,H,\theta,\eta,\psi$ and $\mathrm{d}W=\omega_1\wedge\omega_2$ as above, and write $\phi=\phi_{a,d}$, $A=A_0=c_0H$ and $u=\psi+w_0$. Choose a smooth positive probability volume form $\mathrm{d}V=q\,\mathrm{d}W$ and put $F=\log(2bH/q)$. Then
\[
 \Omega=A\omega_1+b\omega_2,\qquad
 \theta+\ddc u=\Omega+d[D],\qquad
 \mathrm{MA}_\theta(u)=\Omega^2=c_0\mathrm{e}^F\mathrm{d}V.
\]
The model is admissible by \cref{prop:logarithmic-product-model}, and $u\sim\phi$ gives $B_\phi(F)=\log c_0$. Nonconstant $b$ cannot be harmonic, so $\ddc\theta=(\ddc b)\wedge\omega_2\ne0$. Every smooth potential perturbation has the same $\ddc\theta$ and is therefore non-Gauduchon.

For a smooth function $h$ on the product, define its partial Laplacians by $\ddc_i h=(\Delta_i h)\omega_i$. With the mass-one Fubini--Study normalization, $\Delta_i$ is the Laplacian of the unit round sphere.

\begin{lemma}\label{lem:sphere-signed-resolvent}
Let $V>0$ be continuous on $\mathbb P^1$, and let $T$ be a finite signed measure such that
\[
 \int_{\mathbb P^1}(\Delta f-Vf)\,\mathrm{d}T=0
 \qquad(f\in C^\infty(\mathbb P^1)).
\]
Then $T=0$.
\end{lemma}
\begin{proof}
Let $H_t$ be the sphere heat operator. On spherical harmonics of degree $\ell$, it acts by $\mathrm{e}^{-t\ell(\ell+1)}$. The harmonic series defines a smooth heat solution for $t>0$; for smooth initial functions, the maximum principle gives positivity and contraction in the uniform norm. Density of smooth functions extends $H_t$ to a positive, strongly continuous contraction on $C(\mathbb P^1)$. Choose $M\geq\max V$ and define
\[
 R_Mg=\int_0^\infty\mathrm{e}^{-Mt}H_tg\,\mathrm{d}t.
\]
Then $\|R_M\|\leq M^{-1}$ and $(M-\Delta)R_Mg=g$ in distributions. For any $g\in C(\mathbb P^1)$, the map
\[
 f\longmapsto R_M((M-V)f-g)
\]
contracts by at most $1-\min V/M<1$. Its continuous fixed point satisfies $\Delta f=Vf+g$ in distributions, with a continuous right-hand side. The functions $f_\varepsilon=H_\varepsilon f$ are smooth and satisfy
\[
 f_\varepsilon\longrightarrow f,\qquad
 \Delta f_\varepsilon=H_\varepsilon(Vf+g)\longrightarrow Vf+g
\]
uniformly. Thus $(\Delta-V)f_\varepsilon\to g$ uniformly. Testing against the finite measure $T$ gives $\int g\,\mathrm{d}T=0$ for every continuous $g$, hence $T=0$.
\end{proof}

\Needspace{14\baselineskip}
\begin{theorem}\label{thm:singular-product-full-adjoint}\index{Adjoint measure!singular product with continuous residual}
The weight and current
\[
 \rho=\frac1{2ab},\qquad
 Q=\rho\Omega=\frac{A}{2ab}\omega_1+\frac1{2a}\omega_2
\]
satisfy $\mathrm{d}Q=0$ and $\int_X\rho\Omega^2=1$. For smooth real $h$, the measure-valued and logarithmic linearizations at $u$ are
\[
 D_u\mathrm{MA}_\theta(h)=2\Omega\wedge\ddc h,\qquad
 L_uh=\frac{2\Omega\wedge\ddc h}{\Omega^2}
        =\frac{\Delta_1h}{A}+\frac{\Delta_2h}{b}.
\]
The only Borel probability satisfying $\int_XL_uh\,\mathrm{d}\mu=0$ for every smooth real $h$ is
\begin{equation}\label{eq:singular-product-full-support}
 \nu=\rho\Omega^2=\frac Aa\,\mathrm{d}W,\qquad
 \mathcal A_\phi(F)=\{\nu\}.
\end{equation}
If $r\geq0$ is measurable, $\int_Xr\Omega^2=1$ and $\ddc(r\Omega)=0$, then $r=\rho$ almost everywhere.
\end{theorem}
\begin{proof}
We first define the weighted current and variation, then prove invariant-measure uniqueness and apply the supporting inequality on the full space $C(X)$.

\emph{The current and operator.} Since $A/b$ depends only on the first curve, $\mathrm{d}((A/b)\omega_1)=0$ distributionally; this follows also by uniform smooth approximation of its coefficient. Thus $Q$ is closed, and $\int\rho\Omega^2=\int A\omega_1/a=1$. The coefficient $\rho$ is smooth and $\Omega$ is continuous, so the product defining $Q$ is well defined. Uniform positivity of $\Omega$ gives $\Omega+t\ddc h>0$ for sufficiently small signed $t$. The bounded-product calculation used in \cref{thm:moving-logarithmic-selection}, now with the smooth perturbation $th$, gives
\[
 \mathrm{MA}_\theta(u+th)
       =\Omega^2+2t\Omega\wedge\ddc h+t^2(\ddc h)^2.
\]
This proves the two linearizations; mixed components of $\ddc h$ have zero wedge with the diagonal $\Omega$. Distributional Stokes gives $\int L_uh\,\mathrm{d}\nu=2\int\ddc h\wedge Q=0$.

\emph{Invariant probabilities.} Choose real spherical harmonics $e_j$ on the second sphere, with $e_0=1$ and $\Delta_2e_j=-\lambda_je_j$. Only the constant mode has eigenvalue zero. For an invariant probability $\mu$, define finite signed measures $T_j$ on the first sphere by
\[
 \int f\,\mathrm{d}T_j
       =\int_X\frac{f(x)e_j(y)}{A(x)}\,\mathrm{d}\mu.
\]
Testing $L_u$ on $f(x)e_j(y)$ gives
\[
 \int\left(\Delta_1f-\lambda_j\frac Ab f\right)\,\mathrm{d}T_j=0.
\]
For $j=0$, testing on spherical harmonics of the first sphere shows that $T_0=C\omega_1$. Their finite span is uniformly dense, because restrictions of polynomials decompose into harmonics and polynomials are dense on the sphere. The first marginal of $\mu$ is therefore $CA\omega_1$, and probability normalization gives $C=1/a$. For every nonconstant mode, \cref{lem:sphere-signed-resolvent} with $V=\lambda_jA/b>0$ gives $T_j=0$. As these are identities of measures, their test functions may then be multiplied by the continuous function $A$. They determine all spherical harmonic moments of $\mu$. Uniform density of harmonics and of finite product tests gives $\mu=(A/a)\omega_1\wedge\omega_2$.

\emph{Supporting probabilities and weights.} Let $\mu\in\mathcal A_\phi(F)$ and, for smooth real $h$, set
\[
 F_t=F+\log\frac{(\Omega+t\ddc h)^2}{\Omega^2}.
\]
This is continuous, and $u+th$ solves the ordinary-type equation with density $\mathrm{e}^{F_t}$ and scalar $c_0$. Thus $B_\phi(F_t)=B_\phi(F)$, by uniqueness of the scalar. Uniform positivity gives $F_t-F=tL_uh+O(t^2)$ in the uniform norm. Apply \eqref{eq:adjoint-probability-set} for both signs of $t$ to obtain $\int L_uh\,\mathrm{d}\mu=0$. The perturbed density need not depend on the first curve alone: the inequality is the one on all of $C(X)$. Invariant uniqueness and nonemptiness of $\mathcal A_\phi(F)$ from \cref{thm:weak-adjoint-probabilities} now prove \eqref{eq:singular-product-full-support}.

Finally, $\Omega^2/\mathrm{d}W$ is bounded above and below by positive constants, so a normalized $r$ as in the statement belongs to $L^1(\mathrm{d}W)$ and $r\Omega$ is an order-zero current. Distributional Stokes makes $r\Omega^2$ invariant. The probability uniqueness just proved gives $r\Omega^2=\rho\Omega^2$, hence $r=\rho$ almost everywhere.
\end{proof}

\Needspace{12\baselineskip}
\begin{proposition}\label{prop:singular-product-adjoint-approximation}\index{Adjoint measure!fixed-divisor approximation}
Let $w_j\in C^\infty(X)$ be real and suppose
\[
 \Omega_j=\eta+\ddc w_j>0,\qquad \Omega_j\longrightarrow\Omega
       \quad\text{uniformly as }(1,1)\text{-forms}.
\]
Put $u_j=\psi+w_j$ and $F_j=\log(\Omega_j^2/(c_0\mathrm{d}V))$. Let $\rho_j>0$ be the smooth residual adjoint weight, normalized by
\[
 \ddc(\rho_j\Omega_j)=0,\qquad \int_X\rho_j\Omega_j^2=1.
\]
Then $F_j\to F$ uniformly, and the entire sequence satisfies
\[
 \rho_j\Omega_j^2\rightharpoonup\nu,\qquad
 \rho_j\Omega_j\rightharpoonup Q,\qquad
 \rho_j\,\mathrm{d}W\rightharpoonup\rho\,\mathrm{d}W.
\]
Such approximations exist with $w_j\to w_0$ uniformly and can depend on both curve variables.
\end{proposition}
\begin{proof}
Existence of the smooth positive weights is \cref{thm:gauduchon}, applied to the residual metrics $\Omega_j$ on the surface. Put $\nu_j=\rho_j\Omega_j^2$ and $L_jh=2\Omega_j\wedge\ddc h/\Omega_j^2$. Distributional Stokes gives $\int L_jh\,\mathrm{d}\nu_j=0$. Uniform convergence to the positive form $\Omega$ supplies a common determinant lower bound and gives $L_jh\to L_uh$ uniformly for every fixed smooth $h$. Every weak cluster point of the probabilities $\nu_j$ is therefore invariant for $L_u$. By \cref{thm:singular-product-full-adjoint}, it equals $\nu$, proving convergence of the whole sequence. The same determinant convergence gives $F_j\to F$ uniformly.

For a smooth real $(1,1)$ test form $\xi$, the ratios $\Omega_j\wedge\xi/\Omega_j^2$ converge uniformly to $\Omega\wedge\xi/\Omega^2$. Consequently,
\[
 \int_X\rho_j\Omega_j\wedge\xi
    =\int_X\frac{\Omega_j\wedge\xi}{\Omega_j^2}\,\mathrm{d}\nu_j
    \longrightarrow\int_X\rho\Omega\wedge\xi.
\]
Similarly, the positive densities $p_j=\Omega_j^2/\mathrm{d}W$ converge uniformly to $p=\Omega^2/\mathrm{d}W$. For every $g\in C(X)$,
\[
 \int_Xg\rho_j\,\mathrm{d}W
    =\int_X(g/p_j)\,\mathrm{d}\nu_j
    \longrightarrow\int_X(g/p)\,\mathrm{d}\nu
    =\int_Xg\rho\,\mathrm{d}W.
\]
These arguments require no pointwise control of the weights.

For existence, heat-smooth $A$ on the first sphere to positive smooth functions $A_j$ of the same mean $a$, converging uniformly to $A$. Green inversion gives smooth mean-zero $w_j(x)$ with $\ddc w_j=(A_j-a)\omega_1$ and $w_j\to w_0$ uniformly. Their residuals converge uniformly to $\Omega$. Adding $\varepsilon_j h(x,y)$ for a fixed smooth function depending on both variables and $\varepsilon_j\to0$ preserves positivity for all sufficiently small coefficients and supplies nonsymmetric examples. Each $u_j$ retains the fixed logarithmic divisor, and its non-pluripolar equation has scalar $c_0$ and logarithmic density $F_j$.
\end{proof}

\begin{example}\label{ex:singular-product-adjoint-boundary}
Take $a=1$, $d=1/4$, $q=1$, and let $t$ be the sphere height
\[
 t=\frac{|Z_0|^2-|Z_1|^2}{|Z_0|^2+|Z_1|^2},\qquad
 b=2+t,\qquad H=A=\tfrac12+|t|.
\]
Here $c_0=1$ and $F=\log(2bA)$ is Lipschitz but not $C^1$. The residual has continuous positive coefficients with the same loss of differentiability. Its supporting probability differs from normalized Monge--Amp\`ere volume:
\[
 \int_X\Omega^2=4,\qquad
 \int_Xt\,\mathrm{d}\nu=0,\qquad
 \int_Xt\,\frac{\Omega^2}{4}=\frac5{24}.
\]
Even for smooth constant $H$, approximation of the pole by globally smooth potentials need not recover the weak adjoint. For $H=a$, so that $w_0=0$ and $\nu=\mathrm{d}W$, put
\[
 \psi_\delta=d\log(\ell+\delta(1-\ell)),\qquad
 P_\delta=\omega_1+\ddc(\psi_\delta/d),\qquad0<\delta\leq1.
\]
The smooth potentials $\psi_\delta$ decrease pointwise to $\psi$, converge in $L^1$ and smoothly off $D$, but the classical adjoint probabilities of $\theta+\ddc\psi_\delta$ converge to
\[
 \frac{a\omega_1+d\delta_p}{a+d}\wedge\omega_2,
 \qquad p=\{Z_0=0\}.
\]
This limit assigns mass $d/(a+d)>0$ to $D$, and the weighted currents also retain a nonzero divisor component.
\end{example}
\Needspace{8\baselineskip}
\begin{proof}
The height $t$ is uniform on $[-1,1]$, so $\int A\omega_1=1$ and $\int tA\omega_1=0$. Since $\Omega^2=2bA\,\mathrm{d}W$, its mass is four and its normalized height moment is
\[
 \frac12\int_{-1}^1 t^2(\tfrac12+|t|)\,\frac{\mathrm{d}t}{2}
       =\frac5{24}.
\]
The weight is $1/(2b)$ and $\nu=A\,\mathrm{d}W$. Moreover, $\ddc\theta=-2t\,\mathrm{d}W$, so the background has no smooth positive Gauduchon potential perturbation.

For the smoothing family, direct differentiation gives
\[
 \frac{P_\delta}{\omega_1}
       =\frac{4\delta}{((1+\delta)+(1-\delta)t)^2}.
\]
It is positive of mass one and converges weakly to $\delta_p$. Thus the positive smooth forms are $\widehat\Omega_\delta=a\omega_1+dP_\delta+b\omega_2$. Their classical weights and probabilities are
\[
 \widehat\rho_\delta=\frac1{2(a+d)b},\qquad
 \widehat\nu_\delta=\frac{a\omega_1+dP_\delta}{a+d}\wedge\omega_2.
\]
The weighted forms are closed, so these are the normalized classical adjoints. Their limits follow from the curve measure convergence. In particular, the weighted-current limit contains $d[D]/(2(a+d)b(p))$. These smoothings fill the pole and do not satisfy the fixed-divisor and uniform-residual hypotheses of \cref{prop:singular-product-adjoint-approximation}.
\end{proof}

